\documentclass[11pt]{article}
% ============================================================================
% preamble.tex — Estilo común de los manuales de práctica SISELA (FIME / UANL)
%
% Reproduce (de forma aproximada) el estilo de los PDF originales
% (~/Descargas/manuales/Manuales/pN.pdf), de los que no existe fuente .tex.
% Compilar con:  pdflatex  (o  latexmk -pdf pN.tex)
% ============================================================================

\usepackage[utf8]{inputenc}
\usepackage[T1]{fontenc}
\usepackage[spanish,es-tabla]{babel}
\usepackage[a4paper,margin=2.5cm,headheight=15pt]{geometry}

% --- Tipografía: serif para el cuerpo, sans para los títulos ---
\usepackage{libertinus}
\usepackage[scaled=0.92]{helvet}   % sans para encabezados
\usepackage{microtype}

% --- Matemáticas ---
\usepackage{amsmath,amssymb}
\usepackage{siunitx}
\sisetup{detect-all, per-mode=symbol}

% --- Tablas y figuras ---
\usepackage{booktabs}
\usepackage{array}
\usepackage{multirow}
\usepackage{caption}
\captionsetup{font=small,labelfont=bf,labelsep=period}
\usepackage{graphicx}
\usepackage{float}

% --- Diagramas ---
\usepackage{tikz}
\usetikzlibrary{shapes.geometric,arrows.meta,positioning,calc}
\tikzset{
  block/.style   = {draw, rectangle, minimum width=2.6cm, minimum height=0.9cm,
                    align=center, font=\footnotesize},
  term/.style    = {draw, rounded corners, minimum width=2.2cm, minimum height=0.8cm,
                    align=center, font=\footnotesize, fill=black!5},
  decision/.style= {draw, diamond, aspect=2, align=center, font=\scriptsize,
                    inner sep=1pt},
  flow/.style    = {-{Latex[length=2mm]}, thick},
}

% --- Código ---
\usepackage{xcolor}
\definecolor{codebg}{gray}{0.96}
\definecolor{codekw}{rgb}{0.13,0.29,0.53}
\definecolor{codecm}{rgb}{0.40,0.45,0.40}
\usepackage{listings}
\lstdefinestyle{sisela}{
  backgroundcolor=\color{codebg}, basicstyle=\ttfamily\footnotesize,
  keywordstyle=\color{codekw}\bfseries, commentstyle=\color{codecm}\itshape,
  stringstyle=\color{codekw}, numbers=left, numberstyle=\tiny\color{gray},
  numbersep=6pt, frame=none, breaklines=true, showstringspaces=false,
  columns=fullflexible, xleftmargin=1.8em, aboveskip=1em, belowskip=1em,
  literate={á}{{\'a}}1 {é}{{\'e}}1 {í}{{\'i}}1 {ó}{{\'o}}1 {ú}{{\'u}}1
           {ñ}{{\~n}}1 {Á}{{\'A}}1 {°}{{\textdegree}}1 {µ}{{$\mu$}}1
           {✓}{{\checkmark}}1,
}
\lstset{style=sisela}
\newcommand{\code}[1]{\texttt{\small #1}}

% --- Listas ---
\usepackage{enumitem}
\setlist{nosep, leftmargin=1.6em}

% --- Encabezado / pie de página ---
\usepackage{fancyhdr}
\usepackage{lastpage}
\pagestyle{fancy}
\fancyhf{}
\fancyhead[L]{\footnotesize\sffamily Sistemas Eléctricos de Aeronaves — Práctica \practnum}
\fancyhead[R]{\footnotesize\sffamily FIME / UANL}
\fancyfoot[C]{\footnotesize P\'agina \thepage\ de \pageref{LastPage}}
\renewcommand{\headrulewidth}{0.4pt}
\renewcommand{\footrulewidth}{0.4pt}

% --- Títulos de sección en sans-serif ---
\usepackage{titlesec}
\titleformat{\section}{\normalfont\Large\sffamily\bfseries}{\thesection}{0.7em}{}
\titleformat{\subsection}{\normalfont\large\sffamily\bfseries}{\thesubsection}{0.6em}{}
\titleformat{\subsubsection}{\normalfont\normalsize\sffamily\bfseries}{\thesubsubsection}{0.5em}{}

% --- hyperref al final ---
\usepackage[hidelinks,pdfusetitle]{hyperref}

% ============================================================================
% Comandos de portada y cajas
% ============================================================================
\newcommand{\practnum}{N}          % lo redefine cada pN.tex
\newcommand{\practtitle}{}         % idem

\newcommand{\portada}{%
  \thispagestyle{empty}
  \begin{center}
    {\large Universidad Autónoma de Nuevo León}\\[2pt]
    {\normalsize Facultad de Ingeniería Mecánica y Eléctrica}\\[2.2cm]
    \rule{\linewidth}{0.4pt}\\[0.4cm]
    {\Huge\sffamily\bfseries Práctica \practnum}\\[0.5cm]
    {\large \practtitle}\\[0.3cm]
    \rule{\linewidth}{0.4pt}\\[2.5cm]
    {\normalsize Laboratorio de Sistemas Electrónicos de Aeronaves\\(SISELA)}\\[1.5cm]
    {\normalsize Ciclo Escolar 2026}\\[3cm]
  \end{center}
}

\newenvironment{resumen}{\par\noindent\rule{\linewidth}{0.4pt}\par\smallskip
  \small\itshape\noindent\textbf{Resumen Ejecutivo} ---\ }{\par\smallskip
  \noindent\rule{\linewidth}{0.4pt}\par}

\newenvironment{nota}[1][Nota]{\par\smallskip\noindent\rule{\linewidth}{0.4pt}\par\smallskip
  \small\itshape\noindent\textbf{#1} ---\ }{\par\smallskip
  \noindent\rule{\linewidth}{0.4pt}\par\smallskip}

\renewcommand{\practnum}{5}
\renewcommand{\practtitle}{Control de Actuadores I: Modulación por Ancho de Pulso (PWM)\\
y Servomotores — Análisis de Señales — Énfasis Aeronáutico}

\begin{document}
\portada

\begin{resumen}
Esta sesión se enfoca en el control de movimiento para superficies de vuelo mediante
Modulación por Ancho de Pulso (PWM). El estudiante configura los periféricos de
temporización del microcontrolador para comandar servomotores estándar, gestiona la
alimentación de potencia independiente y la lógica de control. La práctica incluye
cuatro modos base (barrido, ángulo manual, pulso directo, potenciómetro) y tres modos
de \textbf{análisis de señales}: medición del \emph{jitter} del PWM con osciloscopio,
demostración de muestreo y \emph{aliasing} inyectando una señal del generador de
funciones en el ADC, y caracterización de la respuesta al escalón del servo-lazo. Se
discute la relevancia de los servomecanismos en sistemas de actuación aeronáuticos
(EMA/EHA, ARINC 429, modos \emph{fail-safe}).
\end{resumen}

\tableofcontents
\newpage

% ============================================================================
\section{Introducción}

El control preciso del movimiento es uno de los pilares de la ingeniería aeronáutica
moderna. Desde el ajuste de las superficies de control de vuelo (alerones, elevadores,
timón, flaps, slats) hasta la actuación de válvulas hidráulicas y el tren de aterrizaje,
se requiere la capacidad de mover componentes a posiciones específicas de forma fiable y
repetible.

La \textbf{Modulación por Ancho de Pulso} (PWM) es la técnica más extendida para lograr
este control con actuadores eléctricos. Un servomotor estándar determina su posición
angular exclusivamente a partir de la duración del pulso activo dentro de un periodo
fijo de \SI{20}{\milli\second} (\SI{50}{\hertz}). El microcontrolador genera esta señal
digital y el circuito de control interno del servo la interpreta para posicionar el eje.

En sistemas \emph{fly-by-wire} (FBW), la cadena de actuación es más compleja: una
computadora de control de vuelo recibe la intención del piloto, calcula las deflexiones
necesarias y transmite comandos digitales de posición a actuadores electro-hidráulicos
o electromecánicos mediante buses como ARINC 429 o AFDX. El principio subyacente es el
mismo: traducir un comando digital en un movimiento mecánico preciso.

Esta práctica añade una dimensión de \textbf{análisis de señales}: el PWM se estudia
como una forma de onda periódica (espectro de armónicos, estabilidad del flanco), el
ADC como un sistema de adquisición (muestreo, aliasing, ENOB) y el servo-lazo como un
sistema dinámico (respuesta al escalón, ancho de banda).

% ============================================================================
\section{Objetivos de Aprendizaje}
\begin{enumerate}
  \item Comprender los principios de PWM: periodo, frecuencia, ciclo de trabajo y ancho
        de pulso, y su relación con la posición de un servomotor.
  \item Describir el funcionamiento interno de un servomotor estándar (motor DC, tren de
        engranajes, potenciómetro de realimentación, circuito de control).
  \item Configurar el periférico PWM del microcontrolador (ESP32 LEDC / RP2040 PWM
        slices) para generar señales estables a \SI{50}{\hertz}.
  \item Calcular los valores de registro (\code{duty\_u16}) necesarios para comandar
        ángulos específicos (0\si{\degree}, 90\si{\degree}, 180\si{\degree}) derivando la
        fórmula de mapeo lineal.
  \item Implementar una clase reutilizable (\code{Servo}) que encapsule la conversión
        ángulo--pulso--registro.
  \item Montar el circuito de alimentación dual: fuente externa para el servo con tierra
        común al MCU.
  \item \textbf{Medir el \emph{jitter} de la señal PWM} con osciloscopio y comparar la
        estabilidad entre plataformas (ESP32 vs RP2040).
  \item \textbf{Demostrar el teorema de Nyquist y el \emph{aliasing}} muestreando con el
        ADC una señal senoidal del generador de funciones y barriendo su frecuencia por
        encima de $F_s/2$.
  \item \textbf{Caracterizar el ADC} (THD, SNR, SINAD, ENOB, SFDR) a partir de una
        captura senoidal.
  \item \textbf{Medir la respuesta al escalón del servo-lazo} (tiempo de subida,
        sobre-oscilación, tiempo de establecimiento, ancho de banda) con realimentación
        de posición.
  \item Contextualizar la actuación PWM dentro de sistemas aeronáuticos reales
        (actuadores EMA/EHA, redundancia, modos \emph{fail-safe}, ARINC 429).
  \item Documentar resultados en formato AIAA.
\end{enumerate}

% ============================================================================
\section{Fundamentos Teóricos}

\subsection{Modulación por Ancho de Pulso}
El PWM varía el tiempo de encendido $t_\text{on}$ de una señal digital dentro de un
periodo constante $T$. El ciclo de trabajo es
\begin{equation}
  D\,(\%) = \frac{t_\text{on}}{T}\times 100,\qquad T = 1/f.
\end{equation}
Para el control de servomotores la frecuencia estándar es $f=\SI{50}{\hertz}$
($T=\SI{20}{\milli\second}$), heredada de los primeros sistemas de radiocontrol. El
servomotor no responde al ciclo de trabajo porcentual sino a la \emph{duración absoluta}
del pulso alto.

\begin{table}[H]\centering
\caption{Relación pulso--ángulo para servomotores estándar.}
\begin{tabular}{cccc}
\toprule
Ancho de pulso & Ángulo & $D$ (\SI{50}{\hertz}) & \code{duty\_u16} \\
\midrule
\SI{1.0}{\milli\second} & 0\si{\degree}   & 5.0\,\%  & 3277 \\
\SI{1.5}{\milli\second} & 90\si{\degree}  & 7.5\,\%  & 4915 \\
\SI{2.0}{\milli\second} & 180\si{\degree} & 10.0\,\% & 6554 \\
\bottomrule
\end{tabular}
\end{table}

\begin{nota}[Rango extendido]
Algunos servos admiten pulsos de \SIrange{0.5}{2.4}{\milli\second} para recorrido
ampliado. El firmware del repositorio usa $t_\text{min}=\SI{500}{\micro\second}$ y
$t_\text{max}=\SI{2400}{\micro\second}$ por defecto, con recorte (\emph{clamp}) para
evitar daño mecánico. Ajustar estos límites experimentalmente para cada modelo.
\end{nota}

\subsection{Servomotores estándar (tipo R/C)}
Un servomotor R/C integra: (1)~un motor DC, (2)~un tren de engranajes reductor,
(3)~un potenciómetro de posición acoplado al eje de salida, y (4)~un circuito de
control (comparador + puente H). El circuito compara el ancho del pulso PWM recibido
con la posición actual del potenciómetro y gira el motor hasta que ambos coinciden:
implementa un \textbf{lazo de control cerrado interno}.

\subsection{Cálculos de mapeo PWM}
El microcontrolador representa el ciclo de trabajo con un registro de 16 bits
(0--65535). Para $T=\SI{20000}{\micro\second}$:
\begin{equation}
  \code{duty\_u16} = \left\lfloor \frac{t_\text{on}\,(\si{\micro\second})}{20000}\times 65535 \right\rfloor.
\end{equation}
La función ángulo $\to$ pulso es una interpolación lineal:
\begin{equation}
  t_\text{on}(\theta) = t_\text{min} + \frac{\theta - \theta_\text{min}}
       {\theta_\text{max} - \theta_\text{min}}\,(t_\text{max} - t_\text{min}).
  \label{eq:map}
\end{equation}

\begin{table}[H]\centering
\caption{Cálculos de temporización (rango \SIrange{500}{2400}{\micro\second}).}
\begin{tabular}{ccccl}
\toprule
$\theta$ (\si{\degree}) & $t_\text{on}$ (\si{\micro\second}) & $D$ (\%) & \code{duty\_u16} & Contexto \\
\midrule
0   & 500  & 2.50  & 1638 & Extremo izquierdo \\
45  & 975  & 4.88  & 3194 & --- \\
90  & 1450 & 7.25  & 4751 & Posición central \\
135 & 1925 & 9.63  & 6308 & --- \\
180 & 2400 & 12.00 & 7864 & Extremo derecho \\
\bottomrule
\end{tabular}
\end{table}

\subsection{Resolución angular}
Con un registro de 16 bits, la resolución teórica es
\begin{equation}
  \Delta\theta = \frac{\theta_\text{max}-\theta_\text{min}}
    {\code{duty\_max}-\code{duty\_min}} = \frac{180\si{\degree}}{7864-1638}
    \approx \SI{0.029}{\degree/bit}.
\end{equation}
En la práctica la resolución la limita la mecánica del servo (juego de engranajes,
histéresis del potenciómetro), típicamente 1--2\si{\degree} para SG90/MG90S.

\subsection{PWM en ESP32 vs RP2040}
\begin{table}[H]\centering
\caption{Comparativa de subsistemas PWM.}
\begin{tabular}{lll}
\toprule
Aspecto & ESP32 (LEDC) & RP2040 (PWM slices) \\
\midrule
Canales & 16 independientes & 8 slices $\times$ 2 = 16 \\
Resolución & configurable (1--20 bits) & fija 16 bits \\
\emph{Jitter} típico & $\sim\SI{100}{\nano\second}$ & $\sim\SI{10}{\nano\second}$ \\
Pin servo (repo) & GPIO18 & GP18 (PWM1 A) \\
Pin ADC (repo) & GPIO34 & GP26 (ADC0) \\
Resolución ADC & 12 bits (0--4095) & 12 bits empacados en 16 (0--65535) \\
Config.\ ADC & \code{atten(11dB)} requerido & sin configuración \\
\bottomrule
\end{tabular}
\end{table}

\begin{nota}[C++/Arduino en RP2040]
El core Arduino para RP2040 usa por defecto \code{analogRead} de \textbf{10 bits}
(0--1023). El firmware C++ de esta práctica llama \code{analogReadResolution(12)} en
\code{setup()} para igualar la resolución de MicroPython.
\end{nota}

\subsection{Contexto aeronáutico: actuadores de control de vuelo}
En aeronaves \emph{fly-by-wire}, las superficies de control no se mueven por cables
mecánicos sino por actuadores que reciben comandos eléctricos digitales:
\begin{enumerate}
  \item \textbf{Actuadores electro-hidráulicos} (EHA): motor eléctrico que acciona una
        bomba hidráulica local (A380, F-35).
  \item \textbf{Actuadores electromecánicos} (EMA): motor eléctrico con reductora; los
        más cercanos al servomotor R/C, pero con potencias de kW, sensores LVDT/RVDT y
        redundancia triple.
  \item \textbf{Actuadores hidráulicos convencionales}: cilindros alimentados por un
        sistema hidráulico central (3000/5000 psi).
\end{enumerate}
La cadena funcional es análoga: comando digital $\to$ conversión a movimiento $\to$
retroalimentación de posición.

\begin{nota}[ARINC 429 Label 101 y 270]
En un sistema de control de vuelo integrado, los comandos de posición viajan como
palabras ARINC 429 de 32 bits: la Label 101 (\emph{Aileron Command}) y la Label 270
(\emph{Flap Command}) transportan el ángulo deseado en formato BNR. El actuador valida
la paridad, verifica el campo SSM y extrae los bits de datos antes de ejecutar el
movimiento.
\end{nota}

\subsection{Modos de falla segura}
En sistemas aeronáuticos certificados, la pérdida de la señal de comando nunca debe
dejar una superficie en posición descontrolada. Los modos \emph{fail-safe} típicos son:
retorno a neutro, congelamiento (\emph{hold last}) y desacoplamiento mecánico. El
firmware del repositorio implementa un concepto similar: al seleccionar \code{q}, el PWM
se desactiva (\code{deinit()}), lo que hace que muchos servos regresen gradualmente a la
posición central por ausencia de señal.

% ============================================================================
\section{Consideraciones de Seguridad y Buenas Prácticas}
\begin{enumerate}
  \item \textbf{Alimentación externa}: nunca alimentar el servo desde los pines 3.3 V o
        5 V del microcontrolador. Los picos de corriente del motor pueden reiniciar el
        MCU, introducir ruido en el ADC o dañar el regulador.
  \item \textbf{Tierra común}: conectar la tierra de la fuente del servo a la del MCU.
        Sin esta referencia común la señal PWM no se interpreta correctamente.
  \item \textbf{Límites mecánicos}: evitar forzar el servo más allá de sus topes físicos.
  \item \textbf{Condensador de desacoplo}: \SIrange{470}{1000}{\micro\farad} electrolítico
        entre VCC y GND del servo, lo más cerca posible del conector.
  \item \textbf{Protección del ADC} (modos de análisis): la salida del generador de
        funciones debe mantenerse dentro de \SIrange{0}{3.3}{\volt} (amplitud
        $\le\SI{2}{Vpp}$, \emph{offset} $+\SI{1.65}{\volt}$). Verificar con el
        osciloscopio antes de conectar. Un pico negativo daña el pin.
\end{enumerate}

% ============================================================================
\section{Equipo y Materiales}
\begin{table}[H]\centering
\caption{Lista de componentes.}
\begin{tabular}{cll}
\toprule
Cant.\ & Componente & Especificación \\
\midrule
1 & Microcontrolador & ESP32 DevKit / Raspberry Pi Pico (RP2040) \\
1 & Servomotor & SG90, MG90S o MG996R (R/C estándar, \SI{50}{\hertz}) \\
1 & Fuente de alimentación & 5 V / 1--2 A, independiente del MCU \\
1 & Potenciómetro & \SI{10}{\kilo\ohm} lineal (modo 4 y realimentación modo 7) \\
1 & Condensador & \SIrange{470}{1000}{\micro\farad} / 10 V electrolítico \\
1 & Protoboard + cables Dupont & --- \\
1 & Generador de funciones & seno hasta $\ge\SI{10}{\kilo\hertz}$, con \emph{offset} DC \\
1 & Osciloscopio & 2 canales, función \emph{Measure} de ancho de pulso \\
--- & PC & Python 3.10+ con \code{tools/sisela\_signal/} \\
\bottomrule
\end{tabular}
\end{table}

% ============================================================================
\section{Etapa de Diseño}

\subsection{Circuito de conexión (alimentación dual)}
\begin{figure}[H]\centering
\begin{tikzpicture}[node distance=1.4cm]
  \node[block] (mcu) {MCU\\(ESP32 / RP2040)};
  \node[block, right=3.2cm of mcu] (servo) {Servomotor\\(SG90 / MG996R)};
  \node[term, above=1cm of servo] (fuente) {Fuente 5 V\\(externa)};
  \draw[flow] (mcu.east) -- node[above,font=\scriptsize]{PWM 50 Hz (señal)} (servo.west);
  \draw[flow] (fuente) -- node[right,font=\scriptsize]{VCC} (servo);
  \draw[thick] (mcu.south) |- ++(0,-0.8) -| (servo.south)
        node[midway,below,font=\scriptsize]{GND común};
  \node[term, below=1.8cm of mcu, font=\scriptsize] (pot) {Pot.\ 10k$\Omega$};
  \draw[dashed,-{Latex}] (pot) -- node[left,font=\scriptsize]{ADC} (mcu.south west);
\end{tikzpicture}
\caption{Esquema de conexión con alimentación dual y potenciómetro (modos 4 y 7).}
\end{figure}

\subsection{Asignación de pines por plataforma}
\begin{table}[H]\centering
\caption{Asignación de pines.}
\begin{tabular}{llll}
\toprule
Señal & ESP32 & RP2040 & Descripción \\
\midrule
PWM servo & GPIO18 & GP18 (PWM1 A) & señal de control 50 Hz \\
VCC servo & fuente ext.\ 5 V & VSYS / ext.\ 5 V & alimentación \\
GND & GND & GND & tierra común \\
ADC (modos 4, 6, 7) & GPIO34 & GP26 (ADC0) & potenciómetro / generador \\
\bottomrule
\end{tabular}
\end{table}

\subsection{Ejemplo numérico}
Para $\theta=90\si{\degree}$, $t_\text{min}=\SI{500}{\micro\second}$,
$t_\text{max}=\SI{2400}{\micro\second}$, usando la ec.~\eqref{eq:map}:
$t_\text{on} = 500 + \frac{90}{180}(2400-500) = \SI{1450}{\micro\second}$;
$\code{duty\_u16} = \lfloor 1450\times 65535/20000\rfloor = 4751$.

% ============================================================================
\section{Procedimiento Experimental}

\subsection{Paso 1: preparación del entorno}
\begin{enumerate}
  \item Instalar MicroPython (o configurar PlatformIO con \code{-DPRACTICE=5}).
  \item Montar el circuito de la figura: servo, fuente externa, tierra común,
        potenciómetro opcional.
  \item Verificar con multímetro: 5 V en VCC del servo, 3.3 V en el GPIO del MCU.
  \item Sincronizar \code{main.py}, \code{boot.py}, \code{lib/servo.py}, \code{lib/siglab.py}.
\end{enumerate}

\subsection{Caso 1: barrido automático 0\si{\degree}--180\si{\degree}--0\si{\degree}}
El modo ejecuta un barrido continuo en pasos de 2\si{\degree} cada \SI{20}{\milli\second}.
Verificar movimiento suave sin vibración excesiva; medir con osciloscopio el ancho de
pulso en los extremos ($\sim\SI{500}{\micro\second}$ y $\sim\SI{2400}{\micro\second}$).

\subsection{Caso 2: control de ángulo manual}
Enviar los ángulos 0, 30, 60, 90, 120, 150, 180. Registrar el pulso reportado
($\si{\micro\second}$) y la posición observada. Medir con osciloscopio el ancho real.

\subsection{Caso 3: calibración por pulso directo}
Enviar valores en microsegundos para encontrar $t_\text{min,real}$ y $t_\text{max,real}$
del servo (donde deja de moverse o hace ruido). Si difieren de 500/2400, actualizar
\code{SERVO\_MIN\_US} / \code{SERVO\_MAX\_US}.

\subsection{Caso 4: control por potenciómetro}
El modo lee el ADC y mapea el valor a un ángulo. Comparar la suavidad entre ESP32
(12 bits) y RP2040 (16 bits) si se dispone de ambas placas.

\subsection{Caso 5: jitter de PWM (osciloscopio)}
\begin{enumerate}
  \item Menú $\to$ 5. Pulso \SI{1500}{\micro\second}, mantener 30 s.
  \item Osciloscopio: CH1 en la señal del servo, \SI{1}{\milli\second\per div}, trigger
        flanco de subida $\sim\SI{1.5}{\volt}$.
  \item \emph{Measure} $\to$ \emph{Pulse Width} con estadística Min/Max/Mean/StdDev.
  \item Exportar una captura larga a CSV y procesar en la PC:
\begin{lstlisting}[language=bash]
python -m sisela_signal characterize --scope --file scopeCH1.csv --col CH1 --mode adc
\end{lstlisting}
  \item Repetir en la otra plataforma y comparar $\sigma$ (esperado: ESP32
        $\sim\SI{100}{\nano\second}$, RP2040 $\sim\SI{10}{\nano\second}$).
\end{enumerate}

\subsection{Caso 6: muestreo y aliasing (generador $\to$ ADC)}
\begin{enumerate}
  \item Generador: seno \SI{1.0}{Vpp}, \emph{offset} \SI{1.65}{\volt},
        \SI{300}{\hertz}. Verificar con el osciloscopio que está dentro de
        \SIrange{0}{3.3}{\volt}.
  \item Conectar a GPIO34 (ESP32) / GP26 (RP2040). Menú $\to$ 6,
        $F_s=\SI{2000}{\hertz}$, $N=2048$.
  \item Guardar el CSV. En la PC:
\begin{lstlisting}[language=bash]
python -m sisela_signal spectrum --file cap.csv --col v --metrics --full-scale 3.3
python -m sisela_signal alias    --file cap.csv --true-f 300
\end{lstlisting}
  \item Barrer la frecuencia del generador: 900, 1000, 1100, 1700, 2100, 3000 Hz. Para
        cada una anotar la frecuencia real y la frecuencia aparente (pico del espectro).
  \item Comprobar que $f_\text{aparente} =
        \lvert ((f + F_s/2)\bmod F_s) - F_s/2\rvert$.
\end{enumerate}

\subsection{Caso 7: respuesta al escalón del servo-lazo}
\begin{enumerate}
  \item Acoplar un potenciómetro al eje del servo (cursor $\to$ ADC, extremos a 3V3 y GND).
  \item Menú $\to$ 7. Ángulos 60 $\to$ 120, $F_s=\SI{500}{\hertz}$, $N=500$.
  \item En la PC:
\begin{lstlisting}[language=bash]
python -m sisela_signal characterize --file step.csv --col v --mode step
\end{lstlisting}
  \item Repetir con escalones 90 $\to$ 95 (pequeño) y 30 $\to$ 150 (grande). ¿Es lineal
        la dinámica?
\end{enumerate}

% ============================================================================
\section{Análisis de Señales}
\label{sec:senales}

\subsection{El PWM como forma de onda periódica}
El tren de pulsos de \SI{50}{\hertz} con ciclo de trabajo $D$ tiene un espectro de
armónicos en $n\cdot\SI{50}{\hertz}$ cuya envolvente es $\lvert\operatorname{sinc}(nD)\rvert$.
El parámetro de calidad crítico para un servo es el \textbf{jitter}: la variación del
ancho de pulso entre periodos, que se traduce en micro-vibración del eje. Se cuantifica
con la desviación estándar $\sigma$ del ancho de pulso (medida en el osciloscopio o
sobre el CSV exportado con \code{sisela\_signal characterize}). La ventaja del RP2040
($\sigma\sim\SI{10}{\nano\second}$ frente a $\sim\SI{100}{\nano\second}$ del ESP32) se
debe a que el ancho de pulso lo fija directamente el \emph{hardware} del PWM slice, sin
intervención del temporizador software.

\subsection{Muestreo, Nyquist y aliasing}
Al muestrear con el ADC a $F_s$ una señal senoidal del generador de frecuencia $f$:
\begin{itemize}
  \item Si $f < F_s/2$ (1\textsuperscript{a} zona de Nyquist), el pico aparece en $f$.
  \item Si $F_s/2 < f < F_s$ (2\textsuperscript{a} zona), aparece un \textbf{alias} en
        $F_s - f$.
  \item Si $f = F_s$, el alias cae en DC (la señal parece ``congelada'').
\end{itemize}
La frecuencia aparente es
$f_a = \lvert ((f + F_s/2)\bmod F_s) - F_s/2\rvert$. El experimento del Caso 6 recorre
varias zonas de Nyquist y verifica esta relación punto por punto.

\subsection{Caracterización del ADC}
A partir de una captura senoidal limpia (coherente si es posible), la FFT permite
calcular (norma IEEE 1241):
\begin{align}
  \text{THD} &= \sqrt{\textstyle\sum_{k\ge 2} P_k / P_1}, &
  \text{SINAD} &= 10\log_{10}\!\frac{P_1}{P_\text{ruido}+P_\text{dist}}, \\
  \text{ENOB} &= \frac{\text{SINAD} - 1.76}{6.02}, &
  \text{SFDR} &= 10\log_{10}\!\frac{P_1}{P_\text{espurio,max}}.
\end{align}
Un ADC de 12 bits nominal presenta un ENOB real de $\sim$9--10 (ESP32 SAR) o $\sim$8.7
(RP2040, según \emph{datasheet}) por efecto del ruido, la no linealidad y el
\emph{jitter} de apertura.

\subsection{Respuesta al escalón del servo-lazo}
El servo es un sistema realimentado (motor + reductora + potenciómetro + comparador).
Con un potenciómetro de realimentación externo se comanda un escalón de ángulo y se
mide la posición real: \textbf{tiempo de subida} (10--90\,\%), \textbf{sobre-oscilación},
\textbf{tiempo de establecimiento} (\(\pm\)\,2\,\%), \textbf{constante de tiempo} $\tau$ y
\textbf{ancho de banda} $f_{-3\,\text{dB}}\approx 1/(2\pi\tau)$ para un modelo de
1\textsuperscript{er} orden. Es el análogo de laboratorio de la dinámica de un actuador
FBW y de los requisitos de ancho de banda de los lazos de control de vuelo.

\subsection{Modo 8 — vista en vivo}
A diferencia de los Modos 5--7 (captura en bloque con \code{BlockSampler}, donde la
precisión de muestreo es crítica y no conviene interrumpirla con impresiones), el
Modo 8 transmite en continuo (\code{stream\_csv}/\code{streamCsv}) el ángulo
comandado y la lectura ADC durante un barrido lento, sin objetivo de medir jitter ni
tiempos críticos. Sirve para observar en tiempo real, mientras el servo se mueve, con
la toolkit compartida:
\begin{lstlisting}[language=bash]
python -m sisela_signal live --port COM5 --menu 8 --cols angle_deg,adc_raw
\end{lstlisting}
Disponible tanto en MicroPython (ambas placas) como en C++ (\code{p5.cpp}), con el
mismo formato de columnas.

% ============================================================================
\section{Alternativa de Trabajo en Casa (Multímetro + Simulación)}
\label{sec:casa}

No todos los casos de esta práctica requieren el mismo equipo. La tabla~\ref{tab:casa5}
resume qué puede hacerse sin osciloscopio ni generador de funciones.

\paragraph{Con multímetro (en casa).} Un multímetro en modo DC promedia la señal PWM:
para un ciclo de trabajo $D$, $V_\text{avg} = \SI{3.3}{\volt}\times D$. Ejemplo: a
90\si{\degree} (\(D=7.5\)\,\%), $V_\text{avg}\approx\SI{0.25}{\volt}$ — medir en el pin de
señal del servo (Caso 2) y comparar contra este cálculo para validar el ancho de pulso
\emph{indirectamente}. Si el multímetro tiene función \textbf{Hz/Duty\,\%}
(común en modelos de gama media), puede leer directamente la frecuencia
($\approx\SI{50}{\hertz}$) y el ciclo de trabajo sin osciloscopio — revisar el manual del
equipo disponible.

\paragraph{Con simulación.} En NI Multisim/Proteus/Wokwi: construir una fuente de pulsos
con los parámetros calculados en la Tabla~2 (periodo \SI{20}{\milli\second}, ancho de
pulso según el ángulo) y un \textbf{osciloscopio virtual} del propio simulador para ``ver''
la forma de onda sin instrumento físico. Es también el lugar ideal para explorar,
\emph{sin riesgo}, qué pasaría con anchos de pulso fuera de rango (p.\ ej. \SI{600}{\micro\second}
o \SI{2600}{\micro\second}) antes de probarlos en un servo real.

\paragraph{Requiere laboratorio presencial (no sustituible).} El \emph{jitter} de PWM
(Caso 5) se mide en nanosegundos: ni un multímetro ni una simulación idealizada (que no
modela las imperfecciones reales del temporizador de hardware) pueden reproducirlo — este
caso exige un osciloscopio real. El muestreo/aliasing (Caso 6) exige una señal senoidal
real y controlada en frecuencia; sin generador de funciones de banco, la alternativa más
cercana es una fuente de audio (teléfono/PC) con divisor resistivo y \emph{offset} DC,
pero requiere precaución adicional para no exceder \SIrange{0}{3.3}{\volt} — se recomienda
dejar este caso para cuando haya acceso al generador de laboratorio.

\begin{table}[H]\centering
\caption{Alternativas de trabajo en casa por caso de estudio.}
\label{tab:casa5}
\begin{tabular}{p{4.6cm}p{3.0cm}p{2.0cm}p{3.2cm}}
\toprule
Caso & Multímetro & Simulación & Requiere laboratorio \\
\midrule
1--4 (barrido, ángulo, pulso, potenciómetro) & Sí ($V_\text{avg}$, Hz/Duty) & Sí & No \\
5 (jitter de PWM) & No & No (idealizada) & \textbf{Sí, osciloscopio} \\
6 (muestreo/aliasing) & Parcial & Parcial & \textbf{Sí, generador} \\
7 (respuesta al escalón) & Sí (indirecto) & Sí & No, 100\,\% viable en casa \\
\bottomrule
\end{tabular}
\end{table}

\begin{nota}[Caso recomendado para trabajo remoto]
El \textbf{Caso 7} (respuesta al escalón) es completamente ejecutable en casa: solo
necesita el servo, un potenciómetro de realimentación y el propio microcontrolador — sin
osciloscopio ni generador. Es el punto de partida sugerido si no se cuenta con acceso al
laboratorio.
\end{nota}

% ============================================================================
\section{Registro de Datos y Análisis}

\begin{table}[H]\centering
\caption{Datos de calibración del servomotor.}
\begin{tabular}{ccccc}
\toprule
$\theta_\text{comando}$ (\si{\degree}) & $t_\text{calc}$ (\si{\micro\second}) &
$t_\text{medido}$ (\si{\micro\second}) & $\theta_\text{observado}$ (\si{\degree}) & Error (\%) \\
\midrule
0   & 500  & & & \\
30  & 817  & & & \\
60  & 1133 & & & \\
90  & 1450 & & & \\
120 & 1767 & & & \\
150 & 2083 & & & \\
180 & 2400 & & & \\
\bottomrule
\end{tabular}
\end{table}

\begin{table}[H]\centering
\caption{Jitter de PWM (Caso 5).}
\begin{tabular}{ccccc}
\toprule
Plataforma & pulso nominal (\si{\micro\second}) & Mean (\si{\micro\second}) &
StdDev (\si{\nano\second}) & pico-a-pico (\si{\nano\second}) \\
\midrule
ESP32  & 1500 & & & \\
RP2040 & 1500 & & & \\
\bottomrule
\end{tabular}
\end{table}

\begin{table}[H]\centering
\caption{Aliasing ($F_s = \SI{2000}{\hertz}$, Caso 6).}
\begin{tabular}{cccc}
\toprule
$f$ generador (Hz) & zona de Nyquist & $f$ aparente medida (Hz) & alias teórico (Hz) \\
\midrule
300  & 1 & & 300 \\
900  & 1 & & 900 \\
1100 & 2 & & 900 \\
1700 & 2 & & 300 \\
2100 & 3 & & 100 \\
3000 & 3 & & 1000 \\
\bottomrule
\end{tabular}
\end{table}

\begin{table}[H]\centering
\caption{Caracterización del ADC (Caso 6) y respuesta al escalón (Caso 7).}
\begin{tabular}{lccccc}
\toprule
Plataforma & THD (\%) & SNR (dB) & SINAD (dB) & ENOB (bits) & SFDR (dB) \\
\midrule
ESP32  & & & & & \\
RP2040 & & & & & \\
\bottomrule
\end{tabular}

\vspace{0.6em}
\begin{tabular}{lcccc}
\toprule
escalón & t.\ subida (ms) & sobre-oscilación (\%) & t.\ estab.\ \(\pm\)\,2\,\% (ms) & $\tau$ (ms) \\
\midrule
60$\to$120 & & & & \\
90$\to$95  & & & & \\
30$\to$150 & & & & \\
\bottomrule
\end{tabular}
\end{table}

\subsection{Gráficas requeridas}
\begin{enumerate}
  \item Ángulo comandado vs ángulo observado (Caso 2) con la recta ideal
        $\theta_\text{obs}=\theta_\text{cmd}$ superpuesta.
  \item Histograma de ancho de pulso del PWM (ESP32 vs RP2040) en el mismo eje.
  \item Espectro de la captura del ADC con el generador a \SI{1100}{\hertz} y
        $F_s=\SI{2000}{\hertz}$, marcando el alias en \SI{900}{\hertz} y la línea de Nyquist.
  \item Frecuencia aparente vs frecuencia real (diagrama de plegado) con los puntos de
        la Tabla de aliasing.
  \item FFT del seno del ADC con marcadores de armónicos (para el ENOB).
  \item Respuesta al escalón del servo-lazo con anotaciones de $t_r$, sobre-oscilación y $t_s$.
\end{enumerate}

% ============================================================================
\section{Análisis de Resultados}
\begin{enumerate}
  \item Linealidad de la actuación: calcular el error porcentual
        $\epsilon = \lvert(\theta_\text{obs}-\theta_\text{cmd})/\theta_\text{cmd}\rvert\times 100$.
  \item Límites reales del servo (Caso 3): comparar $t_\text{min,real}$ / $t_\text{max,real}$
        con los nominales 500/2400 \si{\micro\second}.
  \item Precisión del osciloscopio: comparar $t_\text{on}$ medido vs calculado para
        $\ge 3$ ángulos.
  \item Jitter: ¿la diferencia ESP32 vs RP2040 es observable en la vibración del eje?
        Relacionarla con la resolución angular (\SI{0.029}{\degree/bit}).
  \item Aliasing: ¿los puntos medidos de $f$ aparente coinciden con la teoría? ¿Qué pasa
        exactamente en $f=F_s/2$ y $f=F_s$?
  \item ENOB: ¿dónde se ``pierden'' los bits respecto al valor nominal de 12? ¿Ruido, no
        linealidad, jitter?
  \item Respuesta al escalón: ¿el servo-lazo es de 1\textsuperscript{er} o 2\textsuperscript{o}
        orden? Justificar con la sobre-oscilación y el ajuste de $\tau$.
\end{enumerate}

% ============================================================================
\section{Preguntas de Discusión y Refuerzo}
\begin{enumerate}
  \item Si el PWM tiene resolución de 16 bits y el rango es \SIrange{500}{2400}{\micro\second},
        ¿cuál es el cambio de ángulo más pequeño que el software puede representar? ¿Es
        ese el límite real del sistema o la mecánica impone una cota mayor?
  \item Explicar por qué la tierra común (GND) entre MCU y servo es indispensable. ¿Qué
        sucede si se omite?
  \item El PWM del servo es de \SI{50}{\hertz}. Si muestrearas esa señal con el ADC a
        \SI{80}{\hertz}, ¿qué frecuencia aparente verías? ¿Y a \SI{100}{\hertz} exactos?
  \item Comparar un servo R/C con un actuador electromecánico (EMA) certificado para
        aviación: potencia, velocidad, precisión, sensores de posición, redundancia,
        costo, certificación (DO-178C / DO-254).
  \item En un sistema \emph{fly-by-wire}, el comando de posición viaja como una palabra
        ARINC 429 de 32 bits. Si la Label 101 codifica el ángulo del alerón en BNR con
        resolución de \SI{0.005}{\degree}, ¿cuántos bits de datos se necesitan para cubrir
        $\pm 30\si{\degree}$?
  \item ¿Cuál es la ventaja del RP2040 (\emph{jitter} PWM $\sim\SI{10}{\nano\second}$)
        frente al ESP32 ($\sim\SI{100}{\nano\second}$) para aplicaciones de actuación de
        precisión?
  \item Con $F_s=\SI{2000}{\hertz}$, un tono de \SI{2100}{\hertz} aparece en
        \SI{100}{\hertz}. ¿Cómo lo distinguirías de un tono real de \SI{100}{\hertz} sin
        cambiar la frecuencia del generador? (Pista: cambia $F_s$.)
  \item Un ADC de 12 bits tiene ENOB $\approx 9$. ¿Dónde se han ``perdido'' los 3 bits?
  \item Si añades un filtro de media móvil al comando de ángulo para ``suavizar'' el
        movimiento, ¿qué le pasa al ancho de banda del actuador?
  \item Diseñar (conceptualmente) un modo \emph{fail-safe}: si el MCU no recibe comandos
        válidos durante 3 s, el servo debe regresar a 90\si{\degree}. Describir el algoritmo.
\end{enumerate}

% ============================================================================
\section{Entregables (Formato AIAA)}
\begin{enumerate}
  \item Código fuente: \code{main.py} + \code{lib/servo.py} + \code{lib/siglab.py}
        (o \code{p5.cpp} + \code{common/siglab.h}).
  \item Tabla de calibración completa (7 ángulos con pulsos calculados, medidos y error).
  \item Capturas de osciloscopio: $\ge 3$ ángulos mostrando $t_\text{on}$ y $T$.
  \item Tabla de jitter (Caso 5) y de aliasing (Caso 6).
  \item Tabla de ENOB (Caso 6) y de respuesta al escalón (Caso 7).
  \item Gráficas: las 6 de la \S9.4.
  \item Análisis: respuestas a $\ge 4$ preguntas de la \S11.
  \item Conclusiones: hallazgos, fuentes de error, mejoras propuestas.
  \item Referencias: $\ge 3$ fuentes.
\end{enumerate}

% ============================================================================
\section{Referencias}
\begin{enumerate}
  \item MicroPython Project. \emph{machine.PWM Documentation}.
  \item Raspberry Pi Ltd. (2023). \emph{RP2040 Datasheet}, Sec.\ 4.5: PWM.
  \item Espressif Systems. \emph{ESP32 Technical Reference Manual}, Cap.\ 14: LED Control (LEDC).
  \item TowerPro. \emph{SG90 Micro Servo Datasheet}.
  \item ARINC. (2004). \emph{ARINC Specification 429 Part 1}: Mark 33 DITS.
  \item RTCA. (2011). \emph{DO-178C}: Software Considerations in Airborne Systems.
  \item Moir, I.\ \& Seabridge, A.\ (2008). \emph{Aircraft Systems}, 3rd ed., Wiley.
  \item Pratt, R.W. \emph{Flight Control Systems --- Practical Issues in Design and
        Implementation}, IET, 2000.
  \item IEEE. \emph{IEEE Std 1241-2010 --- Test Methods for Analog-to-Digital Converters}.
\end{enumerate}

% ============================================================================
\appendix
\section{Código Fuente del Repositorio}
\begin{table}[H]\centering
\caption{Archivos relevantes de la Práctica 5.}
\small
\begin{tabular}{p{7.2cm}p{6.0cm}}
\toprule
Ruta & Descripción \\
\midrule
\code{MicroPython/\{ESP32,RP2040\}/P5/main.py} & Programa principal (7 modos) \\
\code{.../P5/lib/servo.py} & Clase \code{Servo} (mapeo ángulo--pulso) \\
\code{.../P5/lib/siglab.py} & Utilidades de análisis de señales en la placa \\
\code{.../P5/docs/analisis\_senales.md} & Guía del módulo de análisis \\
\code{.../P5/tools/servo\_cli.py} & CLI de control remoto por serie \\
\code{C++/SISELA-CPP/src/practices/p5.cpp} & Implementación C++ (7 modos) \\
\code{C++/SISELA-CPP/src/common/\allowbreak\{flight\_controls,siglab\}.h} & Clases comunes \\
\code{tools/sisela\_signal/} & Toolkit PC de análisis de señales \\
\bottomrule
\end{tabular}
\end{table}

\section{Contexto Aeronáutico Complementario}
\textbf{Evolución de los sistemas de actuación.} (1)~Mecánica directa: cables y poleas
(Cessna 172). (2)~Hidráulica asistida: el piloto mueve válvulas, la fuerza la proveen
cilindros hidráulicos (Boeing 737 Classic, A300). (3)~\emph{Fly-by-wire}: señales
eléctricas reemplazan la conexión mecánica; la computadora de control de vuelo (FCC)
interpreta la intención del piloto y envía comandos a los actuadores (A320+, B777+,
F-16, F-35). En FBW la redundancia es mandatoria: un alerón típico tiene dos actuadores
servidos por canales hidráulicos independientes, cada uno comandado por una FCC diferente
con software disimilar (para evitar fallas de causa común, según DO-178C).

\end{document}
